% !TEX program = pdflatex
\documentclass[11pt,a4paper]{article}

\usepackage[utf8]{inputenc}
\usepackage[T1]{fontenc}
\usepackage[spanish,es-nodecimaldot]{babel}
\usepackage{amsmath,amssymb,amsthm,mathtools,mathrsfs}
\usepackage{geometry}
\usepackage{enumitem}
\usepackage[expansion=false]{microtype}
\usepackage{hyperref}
\usepackage{cleveref}
\usepackage{xcolor}
\geometry{margin=2.7cm}
\setlength{\parindent}{0pt}
\setlength{\parskip}{0.6em}

\hypersetup{
	colorlinks=true,
	linkcolor=blue!60!black,
	urlcolor=blue!60!black,
	citecolor=blue!60!black
}

%------------------------------------------------
% Entornos
%------------------------------------------------
\theoremstyle{definition}
\newtheorem{definicion}{Definición}[section]
\newtheorem{ejemplo}{Ejemplo}[section]

\theoremstyle{plain}
\newtheorem{proposicion}{Proposición}[section]
\newtheorem{teorema}{Teorema}[section]
\newtheorem{corolario}{Corolario}[section]
\newtheorem{lema}{Lema}[section]

\theoremstyle{remark}
\newtheorem{observacion}{Observación}[section]

%------------------------------------------------
% Comandos
%------------------------------------------------
\newcommand{\C}{\mathbb C}
\newcommand{\R}{\mathbb R}
\newcommand{\Z}{\mathbb Z}
\newcommand{\N}{\mathbb N}
\newcommand{\I}{[0,1]}
\newcommand{\Hol}{\mathcal H}
\newcommand{\OO}{\mathcal O}
\newcommand{\id}{\operatorname{id}}
\newcommand{\restr}[2]{\left.#1\right|_{#2}}

\title{\textbf{Notas de EDO 2}}
\author{}
\date{}

\begin{document}
	\maketitle

	\section{Miércoles 28 de octubre}

	\begin{definicion}
		Sea $\mathbb{E}$ un espacio vectorial de dimensión finita y $\varphi : \mathbb{R} \times \mathbb{E} \to \mathbb{E}$ el flujo de un campo vectorial $X \in \mathcal{X}^1(\mathbb{E})$. Decimos que:
		\begin{itemize}
			\item $\varphi$ es un \textbf{flujo lineal} si para todo $t \in \mathbb{R}$ la aplicación $\varphi_t : \mathbb{E} \to \mathbb{E}$, $\varphi_t(x) \coloneqq \varphi(t,x)$, es lineal.
			\item $\varphi$ es un \textbf{flujo cuasi-lineal} (o polinomial) si para todo $t \in \mathbb{R}$ la aplicación $\varphi_t : \mathbb{E} \to \mathbb{E}$ es polinomial.
		\end{itemize}
	\end{definicion}

	\begin{proposicion}
		Sea $\mathbb{E}$ un espacio vectorial de dimensión finita y $\varphi : \mathbb{R} \times \mathbb{E} \to \mathbb{E}$ el flujo de un campo vectorial $X \in \mathcal{X}^1(\mathbb{E})$. Entonces $\varphi$ es un flujo lineal si y sólo si existe una única matriz $A \in \mathcal{M}_n(\mathbb{R})$ tal que
		\[
			\varphi_t(x) = e^{tA}\,x, \quad \forall \; t \in \mathbb{R}, \quad \forall \; x \in \mathbb{E}.
		\]
	\end{proposicion}

	\begin{proof}
		\textbf{($\Leftarrow$)} Si $\varphi_t(x) = e^{tA}x$ para alguna matriz $A$, entonces para cada $t \in \mathbb{R}$ la aplicación $\varphi_t$ es la multiplicación por la matriz $e^{tA}$, que es claramente lineal. Luego $\varphi$ es un flujo lineal.

		\textbf{($\Rightarrow$)} Supongamos que $\varphi$ es un flujo lineal, es decir, $\varphi_t \in \operatorname{GL}(\mathbb{E})$ para todo $t \in \mathbb{R}$. Fijada una base de $\mathbb{E}$, identificamos $\mathbb{E} \cong \mathbb{R}^n$ y $\varphi_t$ con una matriz en $\mathcal{M}_n(\mathbb{R})$; abusando de notación escribimos $\varphi_t$ para dicha matriz.

		Como $\varphi$ es un flujo, satisface las propiedades de grupo:
		\begin{enumerate}[label=\roman*)]
			\item $\varphi_0 = \operatorname{id}_{\mathbb{E}}$, es decir, $\varphi_0 = I_n$.
			\item $\varphi_{t+s} = \varphi_t \circ \varphi_s$ para todo $t, s \in \mathbb{R}$.
			\item La aplicación $t \mapsto \varphi_t$ es de clase $C^1$.
		\end{enumerate}
		Diferenciando la identidad de flujo $\varphi_{t+s} = \varphi_t \circ \varphi_s$ con respecto a $s$ en $s = 0$ obtenemos
		\[
			\frac{d}{dt}\varphi_t = \frac{d}{ds}\Big|_{s=0} \varphi_{t+s} = \frac{d}{ds}\Big|_{s=0} \bigl(\varphi_t \circ \varphi_s\bigr) = \varphi_t \cdot \underbrace{\frac{d}{ds}\Big|_{s=0} \varphi_s}_{=:\, A},
		\]
		donde $A \coloneqq \dot{\varphi}_0 \in \mathcal{M}_n(\mathbb{R})$. Así, $t \mapsto \varphi_t$ satisface el problema de valor inicial matricial
		\[
			\left|
			\begin{array}{lcl}
				\dot{M}(t) &=& M(t)\,A, \\
				M(0) &=& I_n.
			\end{array}
			\right.
		\]
		La única solución de este P.V.I. es $M(t) = e^{tA}$ (exponencial matricial), de donde
		\[
			\varphi_t(x) = e^{tA}\,x, \quad \forall \; t \in \mathbb{R}, \quad \forall \; x \in \mathbb{E}.
		\]
		La unicidad de $A$ se sigue de que $A = \dot{\varphi}_0$ queda determinada unívocamente por el flujo. En efecto, si existiese $B \in \mathcal{M}_n(\mathbb{R})$ tal que $e^{tB} = e^{tA}$ para todo $t \in \mathbb{R}$, diferenciando en $t = 0$ se obtiene $B = A$.
	\end{proof}

	\begin{proposicion}
		Sea $\mathbb{E} \cong \mathbb{R}^n$ y $\varphi : \mathbb{R} \times \mathbb{E} \to \mathbb{E}$ un flujo polinomial. Entonces para todo $t \in \mathbb{R}$ existe una única constante $c(t) > 0$ tal que
		\[
			\det\bigl(D\varphi_t(x)\bigr) = c(t), \quad \forall \; x \in \mathbb{E}.
		\]
	\end{proposicion}

	\begin{proof}
		Fijemos $t \in \mathbb{R}$. Dado que $\varphi$ es un flujo polinomial, la aplicación $\varphi_t : \mathbb{E} \to \mathbb{E}$ es polinomial y, siendo parte de un flujo, es un difeomorfismo con inversa $\varphi_{-t}$, que también es polinomial.

		\textbf{Paso 1: $\det(D\varphi_t(\,\cdot\,))$ es una función polinomial.}

		Como $\varphi_t$ es polinomial, su diferencial $D\varphi_t(x)$ es una matriz cuyas entradas son polinomios en $x$. Por definición del determinante (como suma de productos de entradas), $\det(D\varphi_t(x))$ es a su vez un polinomio en $x$. Denotemos $p(x) \coloneqq \det(D\varphi_t(x)) \in \mathbb{R}[x_1, \dots, x_n]$.

		\textbf{Paso 2: $p$ es una constante no nula.}

		De la propiedad de grupo del flujo $\varphi_t \circ \varphi_{-t} = \operatorname{id}_{\mathbb{E}}$, la regla de la cadena da
		\[
			D\varphi_t\!\left(\varphi_{-t}(x)\right) \cdot D\varphi_{-t}(x) = I_n, \quad \forall \; x \in \mathbb{E},
		\]
		de donde, tomando determinantes,
		\[
			\det\!\left(D\varphi_t(\varphi_{-t}(x))\right) \cdot \det\!\left(D\varphi_{-t}(x)\right) = 1, \quad \forall \; x \in \mathbb{E}.
		\]
		Sea $q(y) \coloneqq \det(D\varphi_{-t}(y)) \in \mathbb{R}[y_1, \dots, y_n]$. Sustituyendo $x \mapsto \varphi_t(x)$ en la identidad anterior:
		\[
			p(x) \cdot q\!\left(\varphi_t(x)\right) = 1, \quad \forall \; x \in \mathbb{E}.
		\]
		Como $\varphi_t$ es polinomial, la composición $r(x) \coloneqq q(\varphi_t(x))$ es un polinomio en $x$, y se tiene
		\[
			p \cdot r = 1 \quad \text{en } \mathbb{R}[x_1, \dots, x_n].
		\]
		Como $\mathbb{R}[x_1, \dots, x_n]$ es un dominio de integridad cuyos únicos elementos invertibles son las constantes no nulas, concluimos que $p$ es una constante, es decir, existe $c(t) \in \mathbb{R}\setminus\{0\}$ tal que
		\[
			\det(D\varphi_t(x)) = c(t), \quad \forall \; x \in \mathbb{E}.
		\]

		\textbf{Paso 3: $c(t) > 0$.}

		En $t = 0$, como $\varphi_0 = \operatorname{id}_{\mathbb{E}}$, se tiene $D\varphi_0(x) = I_n$ para todo $x$, luego $c(0) = \det(I_n) = 1 > 0$. La función $t \mapsto c(t) = \det(D\varphi_t(x_0))$ (evaluada en cualquier $x_0$ fijo) es continua en $t$ y satisface $c(0) = 1 > 0$, $c(t) \neq 0$ para todo $t$. Por el teorema del valor intermedio, $c(t) > 0$ para todo $t \in \mathbb{R}$.

		\textbf{Unicidad.} Si existiera $\tilde{c}(t) > 0$ con $\det(D\varphi_t(x)) = \tilde{c}(t)$ para todo $x$, evaluando en cualquier $x_0 \in \mathbb{E}$ se obtiene $\tilde{c}(t) = \det(D\varphi_t(x_0)) = c(t)$.
	\end{proof}

	\textbf{Ejercicios}
	\begin{enumerate}
		\item Si $\varphi$ es un flujo polinomial, demuestre que para cada $t>0$, existe $m=m(t)$ con $gra(\varphi_t)\leq m, \forall t\in [-t,t]$
		\item Si $\varphi$ es polinomial y $X$ es su campo entonces 
		\begin{enumerate}
			\item $X$ es polinomial
			\item $div X$ es constante, digamos $div X = \lambda$ y $det(D\varphi_t(x)= e^{\lambda t})$
		\end{enumerate}
		\item Si $\varphi$ es un flujo polinomial, demuestre que existe $m\in N$ independiente de $t\in R$ tal que
		
		$$\varphi_t(x)=\sum_{|\alpha|\leq m}a_\alpha (t) x^\alpha, a_\alpha(t)\in R^n$$
		
		[1]H. BOSS, G. Meisters, Polynomial flows in the plane, advances in mathematics 55 (1985) 
	\end{enumerate}
	
	\textbf{Ejemplos}
	
	\begin{enumerate}
		\item $X(x)= Ax, A\in \R^{n\times n}$
		Entonces su flujo es $\varphi_t(x)= e^{tA}x$ y su diferencial $Df_t(x)= e^{tA}$. Su determinante es $det(D\varphi_t)(x)= e^{traz(A)}$
	\end{enumerate}
	
	\begin{proposicion}
		Sean $A,B,C\in \mathcal{M}_n(\mathbb{R})$ tales que $BC=CA$. Entonces $e^{tB}C=Ce^{tA}$ para todo $t\in\mathbb{R}$.

		En particular, si $AB=BA$ entonces $e^{tA}B=Be^{tA}$ y
		\[
			e^{tA}e^{tB} = e^{t(A+B)}, \quad \forall\; t\in\mathbb{R}.
		\]
	\end{proposicion}

	\begin{observacion}
		La condición $AB = BA$ es esencial: en general, $e^{A+B} \neq e^A e^B$. La fórmula de Baker-Campbell-Hausdorff cuantifica precisamente la diferencia en términos del conmutador $[A,B] \coloneqq AB - BA$.
	\end{observacion}

	\begin{teorema}[Baker-Campbell-Hausdorff]
		Sean $A, B \in \mathcal{M}_n(\mathbb{R})$. Entonces
		\[
			e^A e^B = e^{A + B + \frac{1}{2}[A,B] + \frac{1}{12}\bigl([A,[A,B]] + [B,[B,A]]\bigr) + \cdots},
		\]
		donde los términos omitidos son combinaciones $\mathbb{Q}$-lineales de conmutadores iterados de $A$ y $B$. Más precisamente, existe un único elemento $Z$ de la álgebra de Lie libre generada por $A$ y $B$ tal que $e^A e^B = e^Z$, dado por la serie
		\[
			Z = \sum_{n=1}^{\infty} Z_n,
		\]
		donde $Z_n$ es una suma de conmutadores anidados de grado $n$ en $A$ y $B$. Los primeros términos son:
		\begin{align*}
			Z_1 &= A + B, \\
			Z_2 &= \frac{1}{2}[A,B], \\
			Z_3 &= \frac{1}{12}\bigl([A,[A,B]] - [B,[A,B]]\bigr), \\
			Z_4 &= -\frac{1}{24}[B,[A,[A,B]]].
		\end{align*}
		La serie converge (en norma operador) cuando $\|A\|$ y $\|B\|$ son suficientemente pequeñas.
	\end{teorema}

	\begin{proof}[Esquema de la demostración]
		Definamos $Z(t)$ implícitamente por $e^{Z(t)} = e^{tA}e^{tB}$. Diferenciando respecto a $t$ y usando la fórmula
		\[
			\frac{d}{dt}e^{Z(t)} = e^{Z(t)}\,\Psi\!\left(\operatorname{ad}_{Z(t)}\right)\!\dot{Z}(t), \quad \Psi(z) \coloneqq \frac{e^z - 1}{z} = \sum_{k=0}^{\infty}\frac{z^k}{(k+1)!},
		\]
		donde $\operatorname{ad}_X(Y) \coloneqq [X,Y]$, se obtiene la ecuación diferencial
		\[
			\dot{Z}(t) = \Psi\!\left(\operatorname{ad}_{Z(t)}\right)^{-1}(A + e^{t\operatorname{ad}_A}B).
		\]
		Resolviendo esta ecuación por iteración con condición inicial $Z(0) = 0$ e integrando en $[0,1]$, se obtiene la serie BCH.
	\end{proof}

	\begin{corolario}
		Si $[A,B] = 0$ (i.e., $A$ y $B$ conmutan), entonces $Z_n = 0$ para todo $n \geq 2$ y la fórmula BCH reduce a
		\[
			e^A e^B = e^{A+B}.
		\]
	\end{corolario}
	Sean $X=Ax$, $Y=Bx$ campos lineales entonces
	$\varphi_t(x)^X=e^{tA}, \quad \varphi_t^Y(x)=e^{tB}$
	
	$$[X,Y]=0\leftrightarrow [A,B]=0$$
	
	$$[X,Y]=\neq 0\leftrightarrow [A,B]\neq 0$$
	
	Sea $X$ un campo vectorial $C^k(k\geq 1)$ definido en un abierto $U\subseteq \R^n$
	
	$$X(f)(p)= X(p)f = df(p)X(p), \forall f\in C^\infty $$
	Entonces $X$ se puede ver como 
	$$X:C^\infty (U) \to C^\infty (U), f\to X(f):U \to \R$$
	es una derivación 
\end{document}
